% LaTeX companion of 05-geometry-of-linear-transformations.typ. Both formats are editable.
\chapter{Geometry of linear transformations}

\section{WS 5, p. 1: examples, scaling, and projection}

在 2-1 我们知道 \(\begin{pmatrix}
0 & {- 1} \\
1 & 0
\end{pmatrix}\) 是 \(\mathbb{R}^{2}\) 中 counter clockwise 转 \(90\) 度的 linear transformation。现在再看几个：

\[A = \begin{pmatrix}
2 & 0 \\
0 & 2
\end{pmatrix},\quad B = \begin{pmatrix}
1 & 0 \\
0 & 0
\end{pmatrix},\quad C = \begin{pmatrix}
{- 1} & 0 \\
0 & 1
\end{pmatrix},\]\[D = \begin{pmatrix}
0 & 1 \\
{- 1} & 0
\end{pmatrix},\quad E = \begin{pmatrix}
1 & 0.2 \\
0 & 1
\end{pmatrix},\quad F = \begin{pmatrix}
1 & {- 1} \\
1 & 1
\end{pmatrix}.\]

The source diagrams apply them to \(\left( {1,2} \right)\) and \(\left( {1, - 1} \right)\):

\begin{itemize}
\tightlist
\item
  \(A\): 放大一倍, sending them to \(\left( {2,4} \right)\) and \(\left( {- 2, - 2} \right)\).
\item
  \(B\): orthogonal projection onto \(x\)-axis, sending both to \(\left( {1,0} \right)\).
\item
  \(C\): reflection about \(y\)-axis.
\item
  \(D\): clockwise 转 \(90\) degrees, sending them to \(\left( {- 2, - 1} \right)\) and \(\left( {1, - 1} \right)\).
\item
  \(E\): 向右 shear, sending them to \(\left( {1.4,2} \right)\) and \(\left( {0.8, - 1} \right)\).
\item
  \(F\): counterclockwise shift + 放大 \(\sqrt{2}\) 倍, sending them to \(\left( {- 1,3} \right)\) and \(\left( {2,0} \right)\).
\end{itemize}

\begin{theorem}{Scalings}

\(\begin{pmatrix}
k & 0 \\
0 & k
\end{pmatrix}\) defines a scaling by \(k\), since \(\begin{pmatrix}
k & 0 \\
0 & k
\end{pmatrix}\begin{pmatrix}
x_{1} \\
x_{2}
\end{pmatrix} = \begin{pmatrix}
{kx_{1}} \\
{kx_{2}}
\end{pmatrix} = k\begin{pmatrix}
x_{1} \\
x_{2}
\end{pmatrix} = kx\).

\end{theorem}

\begin{definition}{\kn{Orthogonal projection}}

Consider a line \(L\) in coordinate plane, 经过 \(\left( {0,0} \right)\). 任何 \(x \in \mathbb{R}^{2}\) 都可以写成 \(x = x^{\parallel} + x^{\perp}\), where \(x^{\parallel}\ \|\ L\) and \(x^{\perp} \perp L\). The transformation \(T(x) = x^{\parallel}\) is the orthogonal projection of \(x\) onto \(L\), denoted \(\text{proj}_{L{(x)}}\).

随意取 \(w\ \|\ L\), \(\text{proj}_{L{(x)}} = \frac{x \cdot w}{w \cdot w}w\). 特别地，如果 \(w\) 是 unit vector \(u = \begin{pmatrix}
u_{1} \\
u_{2}
\end{pmatrix}\ \|\ L\), then \(\text{proj}_{L{(x)}} = \left( {x \cdot u} \right)u\). 这一个 transformation 是 linear 的；with matrix \(P = \begin{pmatrix}
u_{1}^{2} & {u_{1}u_{2}} \\
{u_{1}u_{2}} & u_{2}^{2}
\end{pmatrix}\).

\end{definition}

\begin{definition}{Reflection}

Consider a line \(L\) in coordinate plane, 过 \(\left( {0,0} \right)\), and write \(x = x^{\parallel} + x^{\perp}\). Then \(T(x) = x^{\parallel} - x^{\perp}\) is reflection of \(x\) about \(L\), denoted \(\text{ref}_{L{(x)}}\). 由 Definition 2.2.1, \(\text{ref}_{L{(x)}} = x^{\parallel} - \left( {x - x^{\parallel}} \right) = 2\text{proj}_{L{(x)}} - x = \left( {2P - I_{2}} \right)x\).

\end{definition}

The matrix of \(T\) has form \(\begin{pmatrix}
a & b \\
b & {- a}
\end{pmatrix}\) where \(a^{2} + b^{2} = 1\): \(S = 2P - I_{2} = \begin{pmatrix}
{2u_{1}^{2} - 1} & {2u_{1}u_{2}} \\
{2u_{1}u_{2}} & {2u_{2}^{2} - 1}
\end{pmatrix} = \begin{pmatrix}
{u_{1}^{2} - u_{2}^{2}} & {2u_{1}u_{2}} \\
{2u_{1}u_{2}} & {u_{2}^{2} - u_{1}^{2}}
\end{pmatrix}\).

\section{WS 5, p. 2: rotations and shearing}

\begin{theorem}{Rotations}

\(T(x) = Ax\) is a counterclockwise rotation in \(\mathbb{R}^{2}\) (rotate by \(\theta\)) iff \(A = \begin{pmatrix}
{\cos(\theta)} & {- \sin(\theta)} \\
{\sin(\theta)} & {\cos(\theta)}
\end{pmatrix} = \begin{pmatrix}
a & {- b} \\
b & a
\end{pmatrix}\), where \(a^{2} + b^{2} = 1\).

\end{theorem}

For a nonzero vector \(x\), write its coordinate in polar form \(\left( {r\cos(\varphi),r\sin(\varphi)} \right)\). Rotation gives \(x' = \left( {r\cos\left( {\varphi + \theta} \right),r\sin\left( {\varphi + \theta} \right)} \right)\), so \(x_{1'} = x\cos(\theta) - y\sin(\theta)\) and \(y' = x\sin(\theta) + y\cos(\theta)\), which gives the displayed matrix.

\begin{theorem}{Rotations combined with a scaling}

将 \(v \in \mathbb{R}^{2}\) counterclockwise rotate \(\theta\) 并放大至 \(r\) 倍. Let \(T\left( {a,b} \right) = \begin{pmatrix}
{r\cos(\theta)} \\
{r\sin(\theta)}
\end{pmatrix}\). Then \(T(x) = Ax\), \(A = \begin{pmatrix}
a & {- b} \\
b & a
\end{pmatrix} = r\begin{pmatrix}
{\cos(\theta)} & {- \sin(\theta)} \\
{\sin(\theta)} & {\cos(\theta)}
\end{pmatrix}\).

\end{theorem}

\section{5. Shearing}

Vertical shear:

\[T\left( \begin{pmatrix}
x_{1} \\
x_{2}
\end{pmatrix} \right) = \begin{pmatrix}
x_{1} \\
{kx_{1} + x_{2}}
\end{pmatrix} = \begin{pmatrix}
1 & 0 \\
k & 1
\end{pmatrix}\begin{pmatrix}
x_{1} \\
x_{2}
\end{pmatrix}.\]

Horizontal shear:

\[T\left( \begin{pmatrix}
x_{1} \\
x_{2}
\end{pmatrix} \right) = \begin{pmatrix}
{x_{1} + kx_{2}} \\
x_{2}
\end{pmatrix} = \begin{pmatrix}
1 & k \\
0 & 1
\end{pmatrix}\begin{pmatrix}
x_{1} \\
x_{2}
\end{pmatrix}.\]

\begin{theorem}{Horizontal and vertical shearing}

Horizontal shearing of slope \(kx_{2}\) has matrix \(\begin{pmatrix}
1 & k \\
0 & 1
\end{pmatrix}\); vertical shearing of slope \(kx_{1}\) has matrix \(\begin{pmatrix}
1 & 0 \\
k & 1
\end{pmatrix}\).

\end{theorem}

The source's concluding table records:

\begin{itemize}
\tightlist
\item
  Scaling by \(k\): \(kI_{2} = \begin{pmatrix}
  k & 0 \\
  0 & k
  \end{pmatrix}\).
\item
  Orthogonal projection onto line \(L\): \(\begin{pmatrix}
  u_{1}^{2} & {u_{1}u_{2}} \\
  {u_{1}u_{2}} & u_{2}^{2}
  \end{pmatrix}\), for \(u\ \|\ L\) and \(\left\| u \right\| = 1\).
\item
  Reflection about line \(L\): \(\begin{pmatrix}
  {2u_{1}^{2} - 1} & {2u_{1}u_{2}} \\
  {2u_{1}u_{2}} & {2u_{2}^{2} - 1}
  \end{pmatrix}\).
\item
  Rotation through angle \(\theta\) (逆时针): \(\begin{pmatrix}
  {\cos(\theta)} & {- \sin(\theta)} \\
  {\sin(\theta)} & {\cos(\theta)}
  \end{pmatrix}\).
\item
  Rotation through \(\theta\) with scaling by \(r\): \(r\begin{pmatrix}
  {\cos(\theta)} & {- \sin(\theta)} \\
  {\sin(\theta)} & {\cos(\theta)}
  \end{pmatrix}\).
\item
  Shear: horizontal \(\begin{pmatrix}
  1 & k \\
  0 & 1
  \end{pmatrix}\); vertical \(\begin{pmatrix}
  1 & 0 \\
  k & 1
  \end{pmatrix}\).
\end{itemize}
